%% -*- coding: utf-8 -*-
%% EXAMPLE CONTENT — ships with a new SciTeX workspace. Replace it with yours.
\begin{abstract}
  \pdfbookmark[1]{Abstract}{abstract}

This manuscript is a worked example distributed with SciTeX. It reports no new
research; it exists so that a new workspace shows a complete project --- data,
analysis, figure, and text --- rather than an empty template, and so that every
number printed in it can be checked by the reader in a few seconds.

We ask a question whose answer has been settled since the 1990s: how far does a
purely linear decision rule get on handwritten digit recognition from very
low-resolution images? We take 500 labelled $8\times8$ greyscale digit
thumbnails --- a stratified subset of the public UCI optical-recognition
dataset, redistributed with the project --- split them 70/30 into training and
held-out sets, and fit a linear support vector machine on the raw 64 pixel
values, with no feature engineering and no tuning.

The classifier labels 131 of 150 held-out images correctly (87.3\% accuracy).
Errors are not spread evenly: nine of the ten digits are recognised at 12/15 or
better, while the digit 8 falls to 7/15 and is most often mistaken for a 1.
That asymmetry, visible directly in the confusion matrix, is the interesting
part of the example, and it is what a linear rule on 64 raw pixels should be
expected to do.

Everything reported here is regenerated by a single seeded script in this
project, so the figures and the accuracy above are reproducible offline and
exactly.

\vspace{1em}
\end{abstract}

%%%% EOF
